\section{The joint local-limit conclusion and its realization}
\label{sec:final-realization}
The endpoint of this article remains the parameter-uniform raw mixed
density local limit \eqref{eq:LLT} for the physical record
$(K_{n,R},N_{n,R},T_{n,R})$, with the weighted positive-denominator
interface of Proposition~\ref{prop:conditioning}. Neither the return
section nor the recorded observables are replaced by independent flights,
a smoothed observation, or an unrelated geometric problem.

The change from $Y_R$ to $Y_R^*$ is explicit, and the periodic records
are those of its actual first return. The joint periodic annihilator is
trivial. The adjacent-increment comparison yields a uniform positive
period discrepancy with logarithmic collision length, and
Theorem~\ref{thm:one-step-coercivity} gives its positive-measure
consequence for controlled regular phases. The phase-reconstruction and
Fredholm hypotheses of Theorem~\ref{thm:phase-resolvent-criterion} state
precisely how these estimates would give an induced resolvent bound.
Their verification on the actual induced strong/weak spaces is still
required. No division by an uncontrolled spectral modulus is implicit.

The collision-space spectral input has a different, now completed use:
Theorem~\ref{thm:exponential-return} proves uniform exponential moments
of the genuine return block. It strengthens finite-record stability to
all finite moments. The common-chart construction in
Proposition~\ref{prop:finite-common-patches} gives an explicit
initial-state remainder for the moving-domain density comparison.
The bounded compensation of Section~\ref{sec:collision-covariance}
then proves a uniformly summable collision Green--Kubo formula. The
exact stopped identity gives the actual Gaussian covariance
$D_R=c_*^{-1}\Gamma_R$, its continuity, and its $L^2$ cohomological
kernel. These conclusions no longer depend on an assumed induced
spectral expansion. The stronger induced-correlation criterion of
Proposition~\ref{prop:covariance-closure} remains available, but is not
used to obtain these Gaussian laws. The actual normalized second
moments now converge by Theorem~\ref{thm:marked-quadratic-moments},
and Theorem~\ref{thm:joint-uniform-nondegeneracy} proves uniform positive
definiteness. Its measurable phase argument removes the section term
before using symplectic area and finite physical covers; periodic
evaluation of the transfer function is not required.

At the physical-clock level, the first-order functional law and its
geometric rate comparison remain valid. The stronger maximum estimate
in Theorem~\ref{thm:conditioned-maximal-block} now controls all unfinished
returns on a growing interval, at logarithmic scale, even after
conditioning of polynomially small positive probability. Corollary
\ref{cor:conditional-path-budget} supplies the corresponding explicit
square-root path error. Thus the unfinished-block premise no longer
needs to be assumed in that range of conditioning probabilities.
Theorem~\ref{thm:actual-return-functional} now supplies the unconditioned
functional Gaussian theorem for the completed physical return record;
Corollary~\ref{cor:physical-functional-gaussian} gives unconditioned
physical-time fluctuations as well. Neither an unconditioned functional
limit nor a path comparison under one event justifies changing an exact
lattice conditioning event. That boundary comparison is a separate
weighted local-density problem, as recorded in
Remark~\ref{rem:conditional-interface-scope}.

For the raw inversion itself, the exact edge subtraction and the local
edge criterion retain every correction term. The all-branch residual
must include central critical words and singular itinerary boundaries.
Lemma~\ref{lem:raw-residual-sum} specifies an absolute derivative-norm
criterion for that sum; no uncontrolled word multiplicity is hidden in
it. The required uniform sum has not been established here. Its growth
constants, together with the actual induced operator estimates, must
satisfy the strict splice inequalities \eqref{eq:splice} or directly
prove \eqref{eq:local-tail}. The weighted insertion class needs the same
verification. A positive conditioning denominator then follows from the
raw LLT on central windows by Proposition~\ref{prop:conditioning}.

The resulting dependency chain is exact: the geometry and periodic
coercivity, the genuine exponential return tail, the common finite-record
charts, and the conditioned unfinished-block estimate are proved;
the continuous collision Green--Kubo formula, actual Gaussian laws,
functional limits, and measurable zero-variance cohomology are also
proved. The marked moment laws and full joint nondegeneracy are proved
as well. Quantitative induced phase reconstruction or a direct complete
complementary-frequency estimate, together with the global raw residual
and weighted bounds, remains necessary to instantiate the full density
theorem. No complete unconditional raw
LLT or full A3 density input is asserted before those premises are
verified. The collision spectral theory is used at its stated smooth-multiplier
scope for killing and for the compensated Gaussian limit. No spectral
theorem about the unbounded induced observables is imported.

\subsection*{Acknowledgment and reproducibility}
AI-assisted analysis contributed to this manuscript. The supplementary
source audit fixes the mathematical baseline, the referee report and
all imported source files; the response describes the revisions to each
numbered request. Finite computations check algebra and explicit orbit
margins and do not replace the analytical proofs. Independent human
specialist review is not claimed by this manuscript.
